\documentclass{article}

\usepackage[margin=0.8in]{geometry}
\usepackage{microtype}
\usepackage{booktabs}
\usepackage{graphicx}
\usepackage{caption}
\usepackage{pdflscape}
\usepackage{float}
\usepackage[hidelinks]{hyperref}

\captionsetup{font=small,labelfont=bf}
\setlength{\parindent}{0pt}
\setlength{\parskip}{6pt}

\title{Supporting Information\\How Can You Think That? Deliberation and the Learning of Opposing Arguments}
\author{Gaurav Sood \and Robert C. Luskin \and James S. Fishkin}
\date{September 2026}

\begin{document}
\maketitle

This supplement reports the expanded descriptive and sensitivity results. All tables are
generated by the analysis pipeline. Estimates use the primary coding unless a table states
otherwise. Confidence intervals use CR2 standard errors with Satterthwaite degrees of freedom,
clustering participants by discussion group and treating controls and the participant without a
recorded group as singleton clusters.

\section{Sample composition}

\begin{table}[H]
\centering
\caption{Composition of recruitment, participant, returnee, and control samples}

\begin{tabular}{lrrrrr}
\toprule
Sample & $N$ & Age & Women (\%) & Catholic (\%) & Degree (\%)\\
\midrule
T1 nonparticipant & 444 & 40.0 (436) & 69.1 (444) & 72.6 (424) & 18.8 (441)\\
T2 participant & 124 & 40.7 (124) & 75.8 (124) & 65.8 (120) & 19.5 (123)\\
T3 nonreturning participant & 31 & 37.9 (31) & 74.2 (31) & 66.7 (30) & 12.9 (31)\\
T3 returning participant & 93 & 41.6 (93) & 76.3 (93) & 65.6 (90) & 21.7 (92)\\
T3 control & 150 & 40.0 (149) & 72.7 (150) & 73.9 (142) & 18.7 (150)\\
\bottomrule
\end{tabular}

\begin{minipage}{0.95\textwidth}
\smallskip\footnotesize
\textit{Note:} Cells contain a mean or percentage followed by the nonmissing denominator in
parentheses. T1 nonparticipants completed the recruitment interview but did not attend the
event. T3 nonreturning participants attended but did not complete the follow-up. T3 controls
were recruited separately.
\end{minipage}
\end{table}

\clearpage
\begin{landscape}
\section{Policy-specific results}

\begin{table}[H]
\centering
\scriptsize
\caption{Mean coded reasons and contrasts by policy and prompt}
\resizebox{\linewidth}{!}{%

\begin{tabular}{llrrrrrr}
\toprule
Policy & Prompt & T2 P & T3 P & T3 C & T2 P $-$ T3 C & T3 P $-$ T3 C & T3 $-$ T2\\
\midrule
All-ability schools & Pro-policy & 1.71 (118) & 1.19 (89) & 1.00 (147) & 0.71 [0.31, 1.12] (265) & 0.19 [-0.11, 0.50] (236) & -0.56 [-1.03, -0.10] (85)\\
All-ability schools & Anti-policy & 1.19 (119) & 0.85 (92) & 0.83 (145) & 0.37 [-0.03, 0.77] (264) & 0.02 [-0.24, 0.28] (237) & -0.43 [-0.91, 0.06] (89)\\
Balanced enrolment & Anti-policy & 1.49 (123) & 1.13 (92) & 0.98 (146) & 0.51 [0.22, 0.79] (269) & 0.15 [-0.13, 0.43] (238) & -0.30 [-0.64, 0.04] (92)\\
Balanced enrolment & Pro-policy & 1.88 (118) & 1.60 (78) & 1.26 (127) & 0.62 [0.19, 1.05] (245) & 0.34 [0.01, 0.67] (205) & -0.20 [-0.61, 0.21] (74)\\
Inclusive classrooms & Pro-policy & 1.37 (124) & 1.10 (92) & 1.01 (147) & 0.36 [0.01, 0.72] (271) & 0.09 [-0.19, 0.37] (239) & -0.30 [-0.69, 0.08] (92)\\
Inclusive classrooms & Anti-policy & 0.64 (121) & 0.49 (93) & 0.57 (150) & 0.07 [-0.18, 0.32] (271) & -0.08 [-0.29, 0.13] (243) & -0.12 [-0.49, 0.25] (90)\\
Combined primary and post-primary schools & Pro-policy & 0.92 (123) & 1.02 (93) & 0.72 (149) & 0.19 [-0.08, 0.47] (272) & 0.30 [0.04, 0.55] (242) & 0.11 [-0.25, 0.47] (92)\\
Combined primary and post-primary schools & Anti-policy & 0.79 (122) & 0.92 (91) & 0.81 (149) & -0.02 [-0.35, 0.31] (271) & 0.12 [-0.15, 0.38] (240) & 0.16 [-0.25, 0.56] (90)\\
\bottomrule
\end{tabular}

}
\begin{minipage}{0.96\linewidth}
\smallskip\footnotesize
\textit{Note:} T2 P and T3 P denote participants at the end of the event and follow-up; T3 C
denotes contemporaneous controls. The first two differences compare distinct samples. T3
minus T2 is a within-returnee change. Contrast cells contain an estimate and 95\% confidence
interval. Prompt labels refer to the policy direction, not to the respondent's own position.
Parentheses contain the sample size for each mean or contrast.
\end{minipage}
\end{table}
\end{landscape}

\section{Response-slot order}

\begin{table}[H]
\centering
\caption{Linear probability model for whether a response slot contains a valid reason}

\begin{tabular}{lrrrrrr}
\toprule
Term & Estimate & CR2 SE & CI low & CI high & Rows & Clusters\\
\midrule
T3 participant & 0.076 & 0.036 & 0.002 & 0.149 & 9656 & 170\\
Response-slot order & -0.130 & 0.006 & -0.143 & -0.118 & 9656 & 170\\
Participant $\times$ response-slot order & -0.024 & 0.010 & -0.044 & -0.005 & 9656 & 170\\
\bottomrule
\end{tabular}

\begin{minipage}{0.95\textwidth}
\smallskip\footnotesize
\textit{Note:} The T3 slot-level model includes policy and prompt-side fixed effects. Response
slot order runs from zero for the first offered slot to four for the fifth. The participant
coefficient is therefore the participant--control difference at the first slot; the interaction
shows how that difference changes across later slots. The model describes response production,
not a causal effect of participation.
\end{minipage}
\end{table}

\section{Coding sensitivity}

\begin{table}[H]
\centering
\caption{Sensitivity of the total-repertoire T3 comparison to coding choices}

\begin{tabular}{lrrr}
\toprule
Outcome construction & T3 difference [95\% CI] & $N$ & Clusters\\
\midrule
Inclusive low-credit coding & 1.00 [-0.83, 2.84] & 185 & 134\\
Count every coded mention & 1.15 [-0.65, 2.96] & 185 & 134\\
Coder 1 & 1.28 [-0.23, 2.78] & 243 & 170\\
Coder 2 & 1.08 [-0.38, 2.55] & 243 & 170\\
\bottomrule
\end{tabular}

\begin{minipage}{0.95\textwidth}
\smallskip\footnotesize
\textit{Note:} Estimates are participant minus control. Inclusive coding retains
question-specific low-credit answers. The mention count credits each distinct code rather than
each occupied response slot. Coder-specific outcomes do not require adjudication and therefore
retain more respondents.
\end{minipage}
\end{table}

\begin{table}[H]
\centering
\small
\caption{Sensitivity of directional balance to coding choices}
\resizebox{\textwidth}{!}{%

\begin{tabular}{lrrrr}
\toprule
Construction & T3 difference [95\% CI] & $N$ (T3) & Paired change [95\% CI] & $N$ (paired)\\
\midrule
Inclusive coding: $b$ & 0.03 [-0.10, 0.15] & 177 & -0.20 [-0.34, -0.06] & 51\\
Inclusive coding: $|b|$ & -0.04 [-0.13, 0.05] & 177 & -0.17 [-0.27, -0.07] & 51\\
Coder 1: $b$ & -0.01 [-0.12, 0.10] & 221 & -0.14 [-0.29, 0.01] & 78\\
Coder 1: $|b|$ & -0.04 [-0.14, 0.05] & 221 & -0.10 [-0.21, 0.02] & 78\\
Coder 2: $b$ & 0.03 [-0.09, 0.14] & 221 & -0.12 [-0.25, 0.01] & 78\\
Coder 2: $|b|$ & -0.02 [-0.11, 0.08] & 221 & -0.07 [-0.16, 0.02] & 78\\
\bottomrule
\end{tabular}

}
\begin{minipage}{0.95\textwidth}
\smallskip\footnotesize
\textit{Note:} Directional balance is $b=(s-o)/(s+o)$ and absolute imbalance is $|b|$.
Balance is undefined for a zero repertoire. T3 differences are participant minus control;
paired changes are T3 minus T2. Inclusive coding retains low-credit responses but still
requires completed adjudication. Coder-specific outcomes retain unadjudicated cases.
\end{minipage}
\end{table}

\clearpage
\section{Religious-community descriptions}

\begin{table}[H]
\centering
\caption{T3 total repertoires by religious community}
\resizebox{0.98\textwidth}{!}{%

\begin{tabular}{lrrrrrr}
\toprule
Community & Control mean & $N$ control & Participant mean & $N$ participant & Difference [95\% CI] & $N$ difference\\
\midrule
Catholic & 7.33 & 85 & 7.77 & 43 & 0.44 [-1.97, 2.84] & 128\\
Protestant & 4.82 & 22 & 7.88 & 26 & 3.07 [0.81, 5.32] & 48\\
\bottomrule
\end{tabular}

}
\begin{minipage}{0.95\textwidth}
\smallskip\footnotesize
\textit{Note:} Differences are participant minus control. Respondents outside the two largest
religious communities and respondents with missing affiliation are omitted. These subgroup
estimates are descriptive, were not selected as separate confirmatory tests, and should not be
read as evidence of effect heterogeneity.
\end{minipage}
\end{table}

\section{Coding rules and reproducibility}

The substantive categories follow the original codebook. Codes for ``other'' remain valid;
codes for nonsensical, missing, and vague responses do not. For three prompts, code three
attributes the opposing position to prejudice or ignorance rather than stating a policy reason.
The primary measure excludes that code; the inclusive sensitivity retains it.

Coder labels are normalized to sorted integer sets before agreement is evaluated. If both
coders provide the same set, that set is final. If they differ and the adjudicator provides a
set, the adjudicated set is final. If they differ without adjudication, the primary respondent-
wave outcome is unresolved. Coder-specific sensitivities preserve those observations.

The numbered R scripts create the analysis data, one bundled results object,
headline CSV files, tables, figures, and manuscript macros. Running \texttt{make check} lints
the code, rebuilds every output, runs the validation tests, and compiles the article and this
supplement.

\end{document}
